\documentclass[11pt,letterpaper]{article}
\usepackage[margin=1in]{geometry}
\usepackage[T1]{fontenc}
\usepackage{lmodern,microtype,amsmath,amssymb,amsthm,mathtools}
\usepackage[hidelinks]{hyperref}
\setlength{\emergencystretch}{2em}
\numberwithin{equation}{section}
\newtheorem{theorem}{Theorem}[section]
\newtheorem{proposition}[theorem]{Proposition}
\newtheorem{lemma}[theorem]{Lemma}
\newtheorem{corollary}[theorem]{Corollary}
\DeclareMathOperator{\Ric}{Ric}
\DeclareMathOperator{\Rm}{Rm}
\title{Chiral Weyl energy and the smooth\\four-dimensional Poincar\'e conjecture}
\author{Edward B. Baker III}
\date{Working manuscript, 26 November 2026}

\begin{document}
\maketitle
\begin{abstract}
Let \(I^+(M)\) be the infimum of one chiral Weyl energy over the positive
Yamabe classes on a smooth homotopy four-sphere \(M\). A universal gap
above \(4\pi^2\) for this infimum recognizes the standard smooth sphere.
For \(I^+(M)<8\pi^2\), a companion working theorem obtains a smooth
positive Bach-flat minimizer on the original oriented manifold. Its proof
uses Bour's regularized flow and a bound on the full trace-free Ricci
energy that forces each nonflat blow-up to have one end. Together with
the finiteness theorem of Chang, Qing and Yang, this leaves finitely
many possible nonstandard types in each band \(I^+(M)\le8\pi^2-\delta\),
\(\delta>0\). Two questions remain for this approach to smooth Poincar\'e.
Does every remaining type admit a positive class below \(8\pi^2\)?
Can a nonstandard type carry an attained minimum above the recognition
gap? We also state independent recognition tests for
immersed heights and marked handle diagrams.
\end{abstract}

\section{The problem on a fixed smooth sphere}

Let \(M\) be a closed, connected, oriented smooth manifold homotopy
equivalent to \(S^4\). Freedman's theorem identifies its homeomorphism
type \cite{Freedman1982}. We ask for a diffeomorphism. In Riemannian
dimension four the Hodge star splits real two-forms into two rank-three
bundles, and the Weyl curvature acts on them by trace-free symmetric
endomorphisms \(W^\pm\) \cite{AHS1978}. For a measurable set \(U\)
and a smooth metric \(g\), put
\[
 E_\pm(g;U)=\int_U|W^\pm_g|_{\rm HS}^2\,dV_g,\qquad
 e(g)=E_+(g;M)=E_-(g;M),\qquad T=4\pi^2 .
\]
We take the curvature operator of the unit sphere to be the identity.
The closed energies agree because their difference is \(12\pi^2\tau(M)=0\).
On a region with boundary we keep both integrals. Each chiral measure is
conformally invariant. In the full four-index tensor convention the closed
Weyl integral is \(8e\).

For the scalar curvature \(R\), set
\begin{equation}\label{eq:yamabe}
 Y([g])=\inf_{0\ne u\in H^1(M)}
 \frac{\int_M(6|du|_g^2+R_gu^2)\,dV_g}{\|u\|_{L^4(g)}^2},
 \qquad I^+(M)=\inf_{Y([g])>0} e(g).
\end{equation}
An empty infimum is \(+\infty\). We call \(I^+(M)<2T\) \emph{entry}.
It is equivalent to the existence of a single positive class below \(2T\).
Pullback of the round metric gives \(I^+(S^4)=0\).

The intrinsic argument has three steps. \emph{Entry} constructs such
a class on the type under consideration. \emph{Attainment} produces a
smooth positive class on that same oriented type at exactly \(I^+(M)\).
\emph{Rigidity} would exclude a nonstandard type carrying that minimum.
The attainment theorem used here is the working result of
\cite{Baker2026BourAttainment}. Sections~\ref{sec:thresholds}
and~\ref{sec:rigidity} identify the remaining existence and rigidity
questions. Either the immersed or the topological test in
Section~\ref{sec:immersion} can settle a type independently.

\section{Positive subcritical entry}\label{sec:thresholds}

Write \(Z=\Ric-Rg/4\), \(D=\int_M|Z|^2\,dV\), and
\(A=(\Ric-Rg/6)/2\). Since \(\chi(M)=2\), Chern--Gauss--Bonnet gives
\begin{equation}\label{eq:schouten}
 16\pi^2=2e+\frac1{24}\int_M R^2\,dV-\frac{D}{2},
 \qquad 4\int_M\sigma_2(g^{-1}A)\,dV=16\pi^2-2e .
\end{equation}
In a positive class choose its volume-one Yamabe minimizer. Its scalar
curvature equals \(Y=Y([g])\le Y_{\rm rd}=8\sqrt6\,\pi\)
\cite{Schoen1984Yamabe,LeeParker1987Yamabe}. Thus
\begin{equation}\label{eq:budgets}
 D=4e+(Y^2-Y_{\rm rd}^2)/12\le4e,\qquad Y^2\ge48(2T-e).
\end{equation}
The last inequality is Gursky's Yamabe estimate, in this normalization
\cite[Lemma 4.5]{Bour2010FourthOrderCurvatureFlows}. At zero Yamabe
constant a scalar-flat representative instead gives
\begin{equation}\label{eq:zero-wall}
 e=2T+\tfrac14\int_M|\Ric|^2\,dV .
\end{equation}
In particular a smooth path on \(M\) which starts positive and stays
strictly below \(2T\) stays positive, by continuity of the first
conformal Laplacian eigenvalue.

\begin{proposition}[Geometric form of entry]\label{prop:entry}
Entry on \(M\) is equivalent to the existence of a smooth metric with
\(R>0\) and \(|Z|^2<R^2/12\) everywhere.
\end{proposition}
\begin{proof}
By \eqref{eq:schouten}, an entering class has \(Y>0\) and positive
integral of \(\sigma_2\). The Main Theorem of Chang--Gursky--Yang gives
a representative with \(R>0\) and pointwise positive \(\sigma_2\)
\cite{ChangGurskyYang2002PositiveRicci}. The formula
\(\sigma_2(g^{-1}A)=R^2/96-|Z|^2/8\) gives the stated inequality.
This also gives \(0<2\Ric<Rg\), since each eigenvalue of the trace-free \(Z\) has absolute value at most \(\sqrt{3/4}\,|Z|\). Conversely, integrating the Schouten inequality in a positive-scalar metric gives \(e<2T\).
\end{proof}

One concrete family separates the two requirements. Let \(H_P^5\) be an
oriented five-dimensional handlebody with one zero-handle, a one-handle for each generator, and a two-handle for
each relator of a balanced presentation \(P\) of the trivial group.
Fix the attaching circles and framings. Its spine is
a simply connected acyclic two-complex, so \(H_P\) is contractible.
Turning its handles upside down shows that the boundary has the same
fundamental group. Duality gives the homology of \(S^4\), hence
\(\Sigma_P=\partial H_P\) is a homotopy sphere. Surgeries on points and
circles in its four-dimensional boundary preserve the existence of positive
scalar curvature \cite[Theorem A]{GromovLawson1980Surgery}. An entry
construction must also control the Weyl integral on that entire boundary, with \(E_++E_-<16\pi^2\). For example, the product disk cap matched to a
hyperbolic circle in \(\mathbb H^2(-1)\times S^2(1)\) already costs at
least \(64\pi^2/3\) in the two signs combined. The calculation and
round-boundary sign test are given in \cite[Positivity of a prescribed
exterior]{Baker2026FixedDescent}. A qualifying construction on one
specified boundary would answer entry there. Universal entry, even the
existence of positive scalar curvature on an arbitrary smooth homotopy
four-sphere, remains open.

\section{Recognition at the lower threshold}\label{sec:recognition}

\begin{proposition}[Recognition]\label{prop:recognition}
If \(Y([g])>0\) and \(e(g)\le T\), then \(M\) is diffeomorphic to \(S^4\).
If \(W^+=0\) on \(M\), the class is conformally round.
\end{proposition}
\begin{proof}
The full Weyl integral is \(8e\). In the strict case the sphere theorem
of Chang--Gursky--Yang gives \(S^4\) or \(\mathbb{RP}^4\). At equality
the remaining types in their classification are the Fubini--Study case
or a manifold covered by \(S^3\times S^1\). Simple connectivity and
\(H_2(M)=0\) exclude these alternatives
\cite[Theorems A and B, Remark 2]{ChangGurskyYang2003}.
If \(W^+=0\), signature gives \(W^-=0\). The developing map of the
locally conformally flat structure is a covering of \(S^4\), by compactness,
and therefore a diffeomorphism \cite{Kuiper1949}.
\end{proof}

\begin{theorem}[Uniform recognition above the threshold]\label{thm:recognition-gap}
There is a number \(\varepsilon_*>0\), independent of smooth type,
such that for every closed smooth homotopy four-sphere \(M\),
\[
               I^+(M)\le T+\varepsilon_*
               \quad\Longrightarrow\quad M\cong_{\rm diff} S^4 .
\]
In particular the same inclusive test applies to the energy of any
one actual positive Yamabe class.
\end{theorem}

\begin{proof}
If nonstandard spheres admit positive classes arbitrarily close to
\(T\), choose such classes on \(M_i\), allowing the type to vary,
with \(e_i\to T\) and physical volume-one Yamabe metrics \(g_i\).
Put \(L=\max\{\lambda_{\max}(W^+),\lambda_{\max}(W^-)\}\) and
\(B=4e/3-\int L^2dV\). Trace freeness gives the pointwise bound
\(L^2\le\tfrac23(|W^+|^2+|W^-|^2)\), so \(B\ge0\).
The strict and equality classifications of Chen and Zhu give
\(Y^2<36\int L^2dV\) on any nonstandard smooth homotopy sphere.
Simple connectivity excludes the nontrivial connected-sum factors,
and the exceptional projective plane and compact simply connected
product of surfaces have nonzero second homology
\cite[Theorems 1.4 and 1.6]{ChenZhu2014Conformal}.
Consequently \eqref{eq:budgets} gives
\[
 D_i+3B_i<8(e_i-T),\qquad Y_i\to y_*=8\sqrt3\,\pi,\qquad
 6\|L_i\|_2\to y_* .
\]

In each class take the volume-one modified Yamabe metric \(q_i\)
with \(R_{q_i}-6L_{q_i}=c_i\le0\)
\cite[Lemma 2.1 and Theorem 1.7]{ChenZhu2014Conformal}.
The constant test gives
\(Y_i\le c_i+6\int L_{q_i}dV_{q_i}\le c_i+6\|L_i\|_2\).
Thus \(c_i\to0\), the variance of \(L_{q_i}\) tends to zero, and
\(R_{q_i}\to y_*\) in \(L^2(q_i)\). The conformally invariant
nonnegative spectral deficit has mass \(B_i\to0\), so the sum of
the two squared Weyl densities tends to \(2T\) in \(L^1(q_i)\).
Gauss--Bonnet and the curvature decomposition now give
\[
 \|\Ric_{q_i}-(y_*/4)q_i\|_2\to0,\qquad
 \big\||\Rm_{q_i}|_{\rm tens}^2-64\pi^2\big\|_1\to0 .
\]
The positive conformal factors admit smooth \(C^2\) approximations
at each index. Take volume-one approximations preserving these limits
with \(R\ge-1\), using \(R_{q_i}\ge c_i\). The ordinary Yamabe
inequality absorbs the \(L^2\) scalar error and supplies a uniform
inhomogeneous Sobolev inequality.

For clarity, the compactness step uses Ricci flow of these smooth
auxiliary metrics. Write \(U=|\Rm|_{\rm tens}\), let
\(\phi'=a\phi^2\), \(\phi(0)=8\pi\), with the dimensional
curvature evolution constant \(a\), and set \(v=(U-\phi)_+\).
Initially \(\|v\|_2\to0\). Entropy monotonicity and the
logarithmic Sobolev inequality give on existing short flows a
uniform Sobolev estimate with potential \(R/4\)
\cite[\S3]{Perelman2002}\cite[Theorem 5.5, arXiv version 4]{YeSobolevFlow}.
Since \(R_+\le C(\phi+v)\), small \(\|v\|_2\) absorbs the potential
into an ordinary Sobolev estimate. Testing
\((\partial_t-\Delta)v\le a(v^2+2\phi v)\) against \(2v\)
then absorbs the cubic term and keeps \(\sup_t\|v\|_2=o(1)\)
on a common short horizon. At a space-time maximum of \(tU\) with
\(U=K\) and \(tK\ge2\), the preceding interval of length \(1/K\)
has curvature at most \(2K\). The curvature gradient estimate
gives \(U\ge K/2\) on a ball of radius \(c/\sqrt K\) at that time.
A cutoff iteration in the ordinary Sobolev inequality bounds its
volume below by \(c'K^{-2}\) for large \(K\). This contradicts
the small tail. The resulting \(U\le C/t\) ensures existence on
a common interval \cite[Theorems 7.1 and 8.1]{HamiltonSingularities}.

Let \(\lambda_0=y_*/4\), \(\lambda'=2\lambda^2\),
\(\lambda(0)=\lambda_0\). In an evolving orthonormal frame
\(P=\Ric-\lambda g\) satisfies
\((\partial_t-\Delta)P=2\Rm(P)+2\lambda P\).
The same Sobolev absorption controls \(\int v|P|^2\), hence
\(\sup_t\|P\|_2=o(1)\). The volume equation
\(V'+4\lambda V=-\int\operatorname{tr}P\,dV\) gives
\(V(t)\to(1-2\lambda_0t)^2\). At fixed small positive time,
bounded curvature and derivatives, ball noncollapse and a diameter
bound from volume packing give a smooth closed limit diffeomorphic
to every sufficiently late source on a subsequence
\cite[Theorem 4.7]{CheegerGromovTaylor1982}\cite{HamiltonCompactness}.
It is positive Einstein with scalar square integral
\(16\lambda(t)^2(1-2\lambda_0t)^2=y_*^2\). By Gauss--Bonnet
and signature its chiral energy is exactly \(T\).
Proposition~\ref{prop:recognition} identifies the standard sphere,
so the round metric pulls back to a zero-energy positive class on each
late source. This contradicts their smooth types. The descent paper
gives the complete smoothing estimates
\cite[section on Ricci flow at the recognition endpoint]{Baker2026FixedDescent}.

This excludes nonstandard positive classes with energy below
\(T+\eta\) for some universal \(\eta>0\). Taking
\(0<\varepsilon_*<\eta\) also covers an infimum equal to
\(T+\varepsilon_*\), by its defining approximation.
\end{proof}

\section{Marked spherical cuts}\label{sec:cut}

An embedded round-product comparison collar supplies two pieces of data
for a cut. Its middle sphere marks the boundary identification, and its
retained buffers permit a Yamabe estimate for the two capped classes.
We measure the energy of each entire side in the original metric,
including the part on which the metric is subsequently changed.

\subsection{Marked topology and the energy inequality}

\begin{proposition}[Finite cut criterion]\label{prop:cut-energy}
Let \(g\) be a smooth metric on a closed oriented homotopy four-sphere \(M\). Cut \(M\) along a finite collection of pairwise disjoint smoothly
embedded three-spheres and attach marked standard balls. Give the
capped pieces smooth metrics and the inherited orientations. Write
\[
 m_j=E_+(g;P_j),\qquad E_j=E_+(\widehat g_j),\qquad
 E=E_+(g)=\sum_j m_j .
\]
Suppose the capped metrics have positive Yamabe constants and
\(E_j\leq m_j+\gamma_j\) with \(\gamma_j\geq0\). If
\(\sum_j E_j\leq2T\), in particular if \(E+\sum_j\gamma_j\leq2T\),
at most one capped piece has \(E_j>T\). The original manifold is
oriented-diffeomorphic to that piece if it exists. If all \(E_j\leq T\), the original manifold is
standard. In particular, if \(M\) is nonstandard, the unique piece
\(s\) with \(E_s>T\) satisfies
\begin{equation}\label{eq:cut-discarded}
 \sum_{j\ne s}m_j\leq E-I^+(M)+\gamma_s .
\end{equation}
\end{proposition}

\begin{proof}
Each middle sphere separates, since otherwise a loop could cross it
once, contrary to \(H_1(M;\mathbb Z)=0\). A cut in a current piece
splits it in two: separation holds in the full manifold, and the
connected two-sided collar provides exactly two adjacent components.
The incidence graph of the pieces and cuts is therefore a tree.
Van Kampen and Mayer--Vietoris along its edges, and along the cap
boundaries, give
\[
 \pi_1(M)\cong *_{j}\pi_1(\widehat M_j),\qquad
 H_2(M;\mathbb Z)\cong\bigoplus_j H_2(\widehat M_j;\mathbb Z).
\]
All the factors vanish. By Poincar\'e duality the capped pieces are
simply connected integral homology four-spheres. Hurewicz gives a
map from \(S^4\) representing a fundamental class of each, and the
homological Whitehead theorem makes each a homotopy sphere.

Each cap is an actual positive class on its capped homotopy sphere.
Proposition~\ref{prop:recognition} therefore identifies every piece of
energy \(E_j\leq T\) with \(S^4\). Two energies strictly above \(T\)
would have sum greater than \(2T\), contradicting the total bound.
If two energies equal \(T\), both pieces are standard and every
remaining energy is zero.

Removing the actual cap balls and reusing the original boundary
identifications recovers \(M\). In a standard capped leaf, Palais's
theorem moves the embedded cap ball to a coordinate hemisphere in
the standard sphere. The closed complement of the cap interior is
therefore a standard ball
\cite[Theorem B]{Palais1960Extending}. The resulting boundary
identification extends over a standard ball, since a diffeomorphism
of \(S^3\) is isotopic to an orthogonal map of the same orientation
sign. Use the isotopy in a collar and the orthogonal map towards
the center \cite[p.~553]{Hatcher1983Smale}. Filling the neighboring
piece by the raw leaf or by the cap thus gives the same oriented
diffeomorphism type, by a diffeomorphism that is the identity on
the neighboring piece outside the gluing collar. Root the tree at
the exceptional piece, or at any piece if all are standard.
Successively removing standard leaves identifies \(M\) with the
capped root. In the nonstandard case its metric is an admissible
competitor for \(I^+(M)\), so
\(I^+(M)\leq E_s\leq m_s+\gamma_s\). This gives
\eqref{eq:cut-discarded}.
\end{proof}

The same finite cut proof works with \(T\) replaced throughout
Proposition~\ref{prop:cut-energy} by \(T+\varepsilon_*\), using
Theorem~\ref{thm:recognition-gap} on each positive cap.
In particular the explicit inclusive total \(2T=8\pi^2\) always
suffices once every cap is positive. The scalar-flat wall for a
joined path on the original type still requires a strict subcritical
bound, as in \eqref{eq:zero-wall}.


For one cut there is a useful direct case. If both capped classes
are positive and both raw energies differ from \(E/2\) by at most
\(a\), then \(E/2+a+\max_j\gamma_j\le T\) recognizes both caps and
hence \(M\). For more uneven divisions, the minimizing inequality
\eqref{eq:cut-discarded} brings even a small discarded mass into play.

\begin{corollary}[One-cut nonsplitting]\label{thm:cut}
On a fixed nonstandard type let positive classes have energies
\(e_i\to I^+(M)<2T\). Suppose a marked sphere has full raw sides
\(P_{i,\pm}\) with positive smooth caps of energies at most
\(E_+(g_i;P_{i,\pm})+\gamma_{i,\pm}\). If
\(e_i+\gamma_{i,+}+\gamma_{i,-}\le2T\) at index \(i\), then
\begin{equation}\label{eq:cut-min-bound}
 \min_{\sigma=\pm} E_+(g_i;P_{i,\sigma})
 \le e_i-I^+(M)+\max_{\sigma=\pm}\gamma_{i,\sigma}.
\end{equation}
\end{corollary}
\begin{proof}
In \eqref{eq:cut-discarded} the energy of the discarded side is
bounded by the minimizing excess plus the surcharge on the surviving
side. This proves the inequality, even when the survivor varies with \(i\).
\end{proof}

A short first-order comparison in Proposition~\ref{prop:short-cut}
constructs positive caps by keeping the joined metric strictly below
\(2T\) along a path on \(M\). The inclusive bound above applies once
the capped classes are known to be positive. For a long buffered product
collar the neck paper gives the individual cap estimates directly,
including the localization loss for unequal buffer lengths
\cite{Baker2026MarkedNecks}.

\subsection{Positive entry from the retained exterior}

There is a separate test when a metric has been constructed on a punctured
smooth type and the Yamabe sign of a closed extension is unknown.
Let \(P\) be the compact complement in \(M\) of the interiors of finitely
many disjoint parametrized standard four-balls. Suppose \(h\) is a
smooth metric on \(P\) equal to \(dt^2+\gamma\) on the marked boundary
collars, with \(\gamma\) the unit round three-sphere metric. Write
\(Q_h(v)=\int_P(6|dv|^2+R_hv^2)\,dV_h\).

\begin{proposition}[Round boundary form]\label{prop:round-boundary-form}
There is a smooth closed extension class on \(M\) of energy
\(E_+(h;P)\). It has positive Yamabe constant exactly when every
nonzero \(v\in H^1(P)\), with boundary traces \(f_j\), satisfies
\[
 Q_h(v)+6\sum_j\langle f_j,\sqrt{1-\Delta_\gamma}\,f_j\rangle>0 ,
\]
where \(-\Delta_\gamma\) is nonnegative and the pairing is between
\(H^{1/2}\) and \(H^{-1/2}\), using \(dV_\gamma\). A sufficient condition is positivity of the lowest
eigenvalue of \(-6\Delta_h+R_h\) with
\(\partial_\nu v+v=0\) on each boundary, where \(\nu\) points out of \(P\).
For nonempty boundary this sufficient condition holds whenever
\(R_h\ge0\): vanishing of the Robin form would force a constant function
with zero trace on the connected exterior. Under either sign test,
\(E_+(h;P)\le4\pi^2\) recognizes \(M\), while
\(E_+(h;P)<8\pi^2\) supplies positive subcritical entry.
\end{proposition}

\begin{proof}
Attach unit half-cylinders pointing out from \(P\). Compactify them
in the given markings using a conformal factor equal to one near
\(P\) and to \(e^{-t}\) for large \(t\). In the coordinate
\(r=e^{-t}\) the compactified metric is exactly flat near each added
point. The caps have zero Weyl curvature. On the cylinder the scalar
curvature is six. A spherical harmonic of degree \(l\) has eigenvalue
\(l(l+2)\) for \(-\Delta_\gamma\); this follows by restricting a
homogeneous harmonic polynomial in four variables to the unit sphere.
The least cylinder energy extending its boundary coefficient \(f_l\) is
\[
 6\inf_{w(0)=f_l}\int_0^\infty
       (|w'|^2+(l+1)^2|w|^2)\,dt
   =6(l+1)|f_l|^2 ,
\]
by completing the square; the minimizer is \(f_l e^{-(l+1)t}\).
Orthogonal summation gives the stated boundary form on \(H^{1/2}\)
traces. To justify the energy spaces, let \(X\) denote the exterior
with cylinders, \(f\) the compactifying factor and \(\widehat g=f^2g_X\).
On compactly supported tests the map \(v\mapsto u=fv\) preserves the
Yamabe forms by conformal covariance \cite[\S2]{LeeParker1987Yamabe}.
Adding \(C\|v\|_{2,\widehat g}^2\) adds
\(C\int_X f^2u^2\,dV_{g_X}\) to the cylinder-side form.
For \(C\) sufficiently large the resulting norms are equivalent to the
respective \(H^1\) norms: on \(P\) the factor is one and the scalar is
bounded, and on the cylinders the scalar is six.
Compactly supported tests are dense on \(X\), and tests avoiding the
finitely many added points are dense on the smooth closed metric by
zero point capacity. Thus the correspondence holds on the completions.

Consequently the full quadratic form is the displayed expression
plus a nonnegative sum of zero-trace energies, strictly positive for
each nonzero zero-trace remainder. A nonpositive boundary-form test
extends by its minimizing cap functions to a nonpositive closed test.
On the smooth closed metric the first eigenvalue is attained, so
strict positivity on all nonzero tests is equivalent to positive
Yamabe sign. Finally the eigenvalues \(l+1\) of the boundary operator
are at least one, giving the Robin test by its Rayleigh quotient.
\end{proof}

For a direct recognition test let \(L_h\) be the largest eigenvalue
of either Weyl block on this exterior. If for every nonzero \(v\in H^1(P)\)
\begin{equation}\label{eq:modified-boundary}
 Q_h(v)-6\int_P L_h v^2\,dV_h
      +6\sum_j\langle f_j,\sqrt{1-\Delta_\gamma}\,f_j\rangle>0,
\end{equation}
then \(M\) is standard. Indeed \(L\) has conformal weight minus two
and vanishes on each cap. The preceding extension argument therefore
tests the closed form with potential \(R-6L\) exactly. Its positive
first eigenvalue and bounded potential give coercivity in \(H^1\),
so its modified Yamabe quotient has a positive infimum by Sobolev
embedding. Chen and Zhu's classification and simple connectivity give
the conclusion \cite[Theorem 1.7]{ChenZhu2014Conformal}.
This test uses only the marked exterior. The descent paper gives explicit
round-ended cylinders satisfying it with unbounded exterior Weyl energy
\cite{Baker2026FixedDescent}.

\paragraph{Transfer to an isometric regluing.}
Let \(W\) be a compact oriented smooth contractible four-manifold with
connected boundary \((Y,\gamma)\). Give \(W\) one smooth metric with a
product germ on this actual face. The reflected identity double and
every double glued by an orientation-preserving \(\gamma\)-isometry
are smooth homotopy spheres. If the identity double has positive
Yamabe sign, the free-boundary scalar form on \(W\) is strictly
positive on every nonzero test, by even reflection. Restriction and
addition then give positive sign on every indicated double. Each has
the same chiral energy, namely the sum of the two half energies in
the orientation of \(W\). They also share \(\int L^2dV\), since
orientation reversal interchanges the Weyl blocks. Thus one reflected
metric satisfying the spectral recognition test of
Theorem~\ref{thm:recognition-gap} recognizes all these marked
regluings, including a boundary map that need not extend over \(W\).
A reflected energy below \(2T\) instead gives entry on every such
double. The construction paper derives the free-trace test for a
finite strip and the signature cost at its actual face
\cite[section on short seams]{Baker2026ImmersedHandles}.


The product requirement can be obtained by approximation on the
given half. Suppose a smooth metric \(h\) on \(W\) has positive scalar
curvature, totally geodesic boundary, and induced boundary metric in
the prescribed class \([\gamma]\) in the actual marking. Write
\(B_h=E_+(h;W)+E_-(h;W)\) and \(G_h=\int_W L_h^2dV_h\).
Conformally match \(\gamma\) with zero normal derivative of the
multiplier. In normal coordinates the tangential metric then differs
from its boundary value by \(O(t^2)\). Cutting off this difference
on a layer of width \(b\to0\) gives an exact product with uniformly
bounded curvature and vanishing change in both curvature integrals.
The free scalar form was coercive by positive scalar curvature and
zero mean curvature. The conformal boundary change preserves its
sign, and the layer changes it by \(o(1)\) in the \(H^1\) form norm.
Each compatible isometric double consequently has positive classes
with \(e\to B_h\) and \(P\to4B_h+6G_h\)
\cite[section on the eta-zero filling]{Baker2026ImmersedHandles}.
Either \(B_h\le T+\varepsilon_*\) or \(4B_h+6G_h\le8T+\eta\)
recognizes every one of these fixed smooth types, including equality
with the approximating classes approaching from above.

For the indicated Mazur half there are smooth positive-scalar-curvature
extensions \(g\) of the exact eta-zero boundary metric. The extra
condition needed on some such extension can be tested entirely on
boundary traces. Let \(\mathcal N_g\) be the outward normal derivative
of the \((-6\Delta_g+R_g)\)-harmonic Dirichlet extension and \(H_g\)
the outward mean curvature in the trace convention. A positive
totally geodesic filling in the prescribed class exists exactly when,
for some extension, \(6\mathcal N_g+2H_g\) has positive quadratic
form on all nonzero traces. The positive first Robin eigenfunction
gives a conformal metric with minimal boundary; boundary deformation
then makes it totally geodesic with the same induced metric.
In the converse direction one fixes the target boundary metric by a
conformal multiplier and solves a positive Dirichlet problem to restore
positive scalar curvature. The sign for one actual extension and the
half-energy budget on the resulting metric remain to be established
\cite[section on filling the marked boundary class]{Baker2026ImmersedHandles}.

\section{Attainment on the original smooth type}\label{sec:attainment}

\begin{theorem}[Fixed-type attainment, \cite{Baker2026BourAttainment}]
\label{thm:attainment}
If \(I=I^+(M)<2T\), there is a smooth positive Bach-flat class on
the oriented \(M\) with \(E_+=E_-=I\). On a nonstandard entered
type, \(I=a_M\), the least energy of a positive Bach-flat class in
the strict subcritical band.
\end{theorem}

\begin{proof}[Argument of the companion theorem]
Fix the smooth type first. Choose \(0<\alpha_j\downarrow0\) and
volume-one positive Yamabe metrics on \(M\) with initial excess
\(\epsilon_j=e_{j,0}-I\in[0,\alpha_j^2]\). For fixed \(j\) suppress
the subscript. Bour's functional in our norms is
\[
 {\cal F}^\alpha=2(1-\alpha)e+\tfrac{\alpha}{2}D
                =2e+\tfrac{\alpha}{24}\int R^2\,dV-16\alpha\pi^2 .
\]
The initial value is at most \(2e_0\) by \eqref{eq:budgets}. For late
\(j\), \(e_0/(1-\alpha)\le2T-\eta_0\) for some fixed \(\eta_0>0\).
In particular Bour's pinching hypothesis is strict. Along its negative gradient flow monotonicity gives
\(e(t)\le e_0/(1-\alpha)\) at every smooth time on the original flow.
By \eqref{eq:zero-wall} and \eqref{eq:budgets}, each slice is in a
positive class with \(Y\ge\sqrt{48\eta_0}\), and hence \(e(t)\ge I\).
On the physical flow metric one therefore has
\begin{equation}\label{eq:original-flow}
 D(t)\le D_0+\frac{4(1-\alpha)}{\alpha}(e_0-e(t))
       \le 4I+\frac{4\epsilon_j}{\alpha_j}
       \le 4I+4\alpha_j<32\pi^2 .
\end{equation}
The last inequality has a fixed positive margin for late \(j\).
Here \(D_0\le4e_0\) was used in the starting Yamabe gauge. The second
formula for \({\cal F}^\alpha\) also gives
\(\int R(t)^2\,dV_{g(t)}\le Y_0^2+48\epsilon_j/\alpha_j\), where
\(Y_0\) is the initial Yamabe constant.

\emph{The two outcomes at fixed \(\alpha\).}
Bour's Theorem~1.11 gives a closed \({\cal F}^\alpha\)-critical
subsequential limit or a complete noncompact nonflat scalar-flat
Bach-flat limit at a curvature maximum
\cite[Theorems 1.4, 1.9 and 1.11]{Bour2010FourthOrderCurvatureFlows}.
The companion justifies the short-time uniqueness and bounded-slab
smoothing used in this strict case \cite{BahuaudHelliwell2015Uniqueness,Baker2026BourAttainment}.
The closed endpoint stays on \(M\) and has constant scalar curvature
by the trace of the regularized equation. It can have nonzero Bach tensor.

For a pointed limit \(X\), compactly supported tests and compact
exhaustion pass the Yamabe lower bound, each chiral upper bound and
\eqref{eq:original-flow} to the limit. The homogeneous Sobolev estimate
and finite curvature give finitely many asymptotically locally Euclidean
ends with differentiated decay
\cite{TianViaclovsky2005ALE,TianViaclovsky2008Volume,AcheViaclovsky2011Obstruction}.
On a connected neighborhood of the \emph{whole} truncated core, including
the buffers of every end, take one oriented comparison embedding into the
same source slice. Its disjoint ends can all be straightened to round
quotient products on the still closed source with arbitrarily small
chiral cost. The spherical-link obstruction at \(E_+<9\pi^2\) forces
each quotient group to be trivial \cite[A quotient obstruction on the
closed source]{Baker2026MarkedNecks}. If there are \(k\) ends, capping
the core gives a homotopy sphere by the marked-cut topology, so
\(\chi(X)=2-k\). The Euler end term can be computed using only the closed formula.
On each dyadic end annulus splice the differentiated ALE metric to the
Euclidean metric. In the radius-one rescaling the change of its total
quadratic Gauss--Bonnet integral is \(o(1)\). Close each resulting flat
outlet of radius \(r\) with a fixed scaled ball whose smooth warped
metric continues the increasing flat radius across the seam and closes
at a round pole. This cap contributes exactly \(16\pi^2\): joined to the
flat radius-\(r\) ball it gives a smooth \(S^4\), and the flat ball
contributes zero. The capped core has Euler characteristic two.
Letting \(r\) tend to infinity in its closed formula gives
\[
 D_X=2\big(E_+(X)+E_-(X)\big)+32\pi^2(k-1).
\]
The strict bound in \eqref{eq:original-flow} forces \(k=1\).

Conformal compactification and removal of its point give a smooth
positive Bach-flat homotopy sphere \(N\). In particular its chiral
energies agree and \(D_X=4e(N)\). The closed critical gap and
nonflatness imply \(e(N)>T\)
\cite[Strict recognition and the closed critical gap]{Baker2026MarkedNecks}.
The smooth type still has to be identified. Remove a small normal ball
in the smooth compactifying chart and embed its complement with both
collar buffers into one original slice. In each sign the upper limit of the whole raw
exterior energy is at most \(2T-\eta_0-e(N)+o(1)<T\).
Choose the buffer length first, then the ball radius, then the comparison
index. The buffered cap estimate gives a positive cap on that actual
exterior of energy below \(T\). Recognition makes this cap a standard
sphere. Disk-neighborhood extension identifies the original exterior
with the deleted normal ball, matching the inverse core map on their
common collar \cite[Theorem C]{Palais1960Extending}. Gluing the two
maps gives an oriented diffeomorphism \(M\cong N\), with smooth
inverse in the same marking. In particular \(I\le e(N)\) and exhaustion
gives \(e(N)\le e_0/(1-\alpha)\).
Keeping the initial Yamabe term in \eqref{eq:original-flow} gives
\(4(e(N)-I)+(Y_{\rm rd}^2-Y_0^2)/12\le4\epsilon_j/\alpha_j\).
Thus pointed outcomes along infinitely many original flows require
their initial Yamabe constants to approach the round value.

\emph{The limit as \(\alpha_j\to0\).}
Retain the volume-one physical representative of a closed regularized
endpoint. For a compactified pointed endpoint take its volume-one
Yamabe representative on the identified \(M\) and recompute
\(D\le4e\) there. The resulting closed metrics \(q_j\) have
\(I\le e(q_j)\le e_{j,0}/(1-\alpha_j)\), uniformly positive Yamabe
constants, bounded constant scalar curvatures and a common strict
Ricci-energy bound below \(32\pi^2\). In the constant-scalar Euler
equation the leading term is \(\tfrac12\nabla^*\nabla Z\), and the
parameter \(p_j=\alpha_j\) or \(0\) enters only in a bounded zeroth-order
coefficient. On curvature-normalized balls the Yamabe inequality gives
noncollapse. In harmonic charts the Ricci equation for the metric and
the elliptic equation for \(Z\) give uniform interior derivative bounds
\cite[The limit of critical metrics]{Baker2026BourAttainment}. These are
stationary estimates uniform in \(j\).

Bounded curvature now gives a closed smooth Bach-flat limit on \(M\)
with energy \(I\). At unbounded curvature a rescaling about a maximum
gives another complete nonflat scalar-flat Bach-flat limit. The same
marked compactification argument applies to this sequence of \(q_j\).
Its positive class on \(M\) satisfies
\(I\le e(N)\le\liminf_j e(q_j)=I\). In this last case the exterior mass in either sign tends to zero when the source index tends to infinity first and the normal ball radius then tends to zero. For a pointed output at a fixed positive parameter the corresponding upper bound is \(e_0/(1-\alpha)-I\). This proves attainment in both
cases. On a nonstandard type the minimum lies above \(T\). Every positive
exact class is admissible for \(I\), and the attaining class is itself
exact, proving \(a_M=I\).
\end{proof}

\section{Finiteness in a fixed band}\label{sec:exact-band}

\begin{corollary}[Finiteness for the infimum]\label{cor:entered-type-attainment}
For each \(\delta>0\), only finitely many oriented nonstandard smooth
homotopy four-sphere types have \(I^+(M)\le 2T-\delta\).
\end{corollary}
\begin{proof}
Even at equality this infimum supplies an entering class, since
\(\delta>0\). Apply Theorem~\ref{thm:attainment} to take a positive
Bach-flat class on the same type at the infimum itself. Its total
Weyl integral is uniformly bounded. In our Schouten normalization,
\[
       \int_M\sigma_2(g^{-1}A)\,dV=(2T-e)/2\ge\delta/2 .
\]
These are the hypotheses of the finiteness theorem of Chang, Qing and
Yang \cite[Theorem B]{ChangQingYang2006Gap}. Their Schouten tensor in
the statement is \(2A\), so their lower bound is \(2\delta\).
Each unoriented smooth type has at most two oriented types.
\end{proof}

Consequently, for any sequence of pairwise distinct nonstandard entered types, their infima must tend to \(2T\) from below. On the positive
Bach-flat classes in a fixed nonstandard band \(e\le2T-\delta\), the spectral deficit
\(36\int L^2\,dV-Y^2\) has a strictly positive attained minimum whenever
the family is nonempty. The least exact value on each type is isolated
among exact values. These refinements use the compactness or strict
energy drop of the marked-neck theorem
\cite[First fixed-type critical value]{Baker2026MarkedNecks}
\cite[The spectral deficit of exact critical classes]{Baker2026FixedDescent}.
Neither finiteness nor isolation excludes an individual nonstandard type.


\section{Rigidity of an attained minimum}\label{sec:rigidity}

Assume now that an entered type is nonstandard, and take the class from
Theorem~\ref{thm:attainment}. Its energy satisfies
\begin{equation}\label{eq:band}
             T+\varepsilon_*<e(g_*)=I^+(M)<2T .
\end{equation}
For any smooth symmetric tensor \(h\), positive Yamabe sign persists
along \(g_*+th\) for small \(|t|\). The first variation of the Weyl
integral thus vanishes, giving the Bach equation
\cite[\S3]{ChangGurskyYang2003}. In particular, its second variation satisfies
\[
                D^2E_+(g_*)[h,h]\ge0 .
\]
At this critical point the Hessian vanishes on infinitesimal conformal changes and
diffeomorphisms. A negative direction transverse to these symmetries
would give a positive competitor of strictly lower energy on \(M\).
A finite marked change with that same property would also suffice.

In a volume-one Yamabe representative the trace-free Ricci tensor \(Z\)
is divergence-free. Put \(C_{ijk}=\nabla_i Z_{jk}-\nabla_j Z_{ik}\).
The Ricci-direction calculation in the descent paper gives the necessary
inequality
\[
 0\le 4D^2E_+(g_*)[Z,Z]
   =\int_M\left(\tfrac7{12}|Z|^4-\operatorname{tr}((Z^\sharp)^4)
                     -3Z^{ia}C_{ijk}C_a{}^{jk}\right)dV .
\]
Its derivation uses constant scalar curvature and Bach-flatness
\cite[The reduced Ricci-flow test]{Baker2026FixedDescent}.
The signed Cotton contraction is one obstruction to obtaining a
contradiction from this test alone.

For a positive Einstein representative, \(Z=0\) and
\(e=2T-Y^2/48\). There is a different test on transverse-traceless
symmetric tensors. Take the Lichnerowicz Laplacian with leading term
\(\nabla^i\nabla_i\). An eigenvalue of \(-\Delta_L\) in \((R/3,R/2)\)
would give a negative Weyl second variation, since its sign is that of
\((R/3-\lambda)(R/2-\lambda)\)
\cite[Theorem 3.10, \(n=4,\ \tau=-1/3\)]{GurskyViaclovsky2015}.
A putative minimizing Einstein metric must avoid this open interval.

On the original positive class, one may also choose a volume-one representative with positive
constant \(\sigma_2\), unique in a nonround admissible conformal class
\cite[Theorem 1.2]{ChangGurskyYang2003}\cite[Theorem 1.5]{GurskyStreets2018TwoYamabe}.
Its positive Ricci curvature and the conformally invariant Bach
equation hold on the same representative. A rigidity argument in
either gauge still needs a sign or a classification beyond
\eqref{eq:band}. A symmetric ansatz on the standard sphere tests
criticality in that class of metrics. The global minimum on the
standard sphere is already zero.

The constant sequence \(g_i=g_*\) explains the limit of the disk
comparisons. A marked disk can carry definite Weyl energy on both sides
at every index of that sequence. If a positive replacement on the same
\(M\) starts at energy \(e_i\), its net saving is at most \(e_i-I^+(M)\).
At \(g_*\) a contradiction requires an actual strict saving. When a
collar is used to infer positivity, one must bound the maximum energy
along the joined path strictly below \(2T\), keep the protected core,
and include each connector at its actual width in the two complete
chiral budgets. The cap tests counting already positive pieces in
Proposition~\ref{prop:cut-energy} include equality. The deformation
and surgery papers measure these different costs on a specified carrier
\cite{Baker2026CollarDeformation,Baker2026CollarProgram}.

\section{An immersion test for entry and recognition}\label{sec:immersion}

There is another way to specify a metric on the given smooth domain.
Every smooth homotopy four-sphere is stably parallelizable
\cite[Theorem 3.1]{KervaireMilnor1963Groups}. In rank five the
stabilization is already in the stable range over a four-dimensional
complex. Hence \(TM\oplus\mathbb R\) is trivial, and the immersion
theorem gives a smooth \(F:M\to\mathbb R^5\)
\cite[\S6]{Hirsch1959Immersions}. Give \(M\) the induced metric
\(g\), write \(A\) for its shape operator, \(H=\operatorname{tr} A\),
\(A^\circ=A-(H/4)I\), and set
\[
 J=\int_M(-\det A)_+\,dV_g,\qquad B=\|A^\circ\|_{L^4(g)}^2 .
\]
Self-intersections are counted on the smooth domain.

\begin{theorem}[Immersed recognition and entry]\label{thm:immersed-entry}
Each of the following conditions on one immersion implies
\(M\cong S^4\) in the smooth category:
\[
       J<8\pi^2/3,\qquad E_+(g)\le64\pi^2/9,\qquad
                     B\le(16\sqrt2/3)\pi .
\]
The conditions are alternatives. The energy condition holds equivalently for any representative
conformal to the induced metric. It requires no Yamabe sign assumption.
On a hypothetical nonstandard type, an immersion with
\(E_+(g)<8\pi^2\) has exactly two index-two critical points in
more than three quarters of all height directions (normalized
spherical measure).
\end{theorem}

\begin{proof}
Choose the unit normal so that its Gauss map \(n:M\to S^4\)
has oriented pullback tangent bundle \(TM\). Euler classes give
\(\deg n=1\), and its signed Jacobian is \(\det A\).
Put \(\omega=|S^4|=8\pi^2/3\).
For almost every \(u\in S^4\), the height \(\langle F,u\rangle\)
is Morse, with critical points \(n^{-1}(\{u,-u\})\). The area
formula gives the Chern--Lashof average
\cite{ChernLashof1957Curvature}
\[
 \frac1{\omega}\int_{S^4}\#\operatorname{Crit}\langle F,u\rangle\,du
  =\frac2{\omega}\int_M |\det A|\,dV_g=2+4J/\omega .
\]
If \(J<\omega\), some Morse height has at most four critical
points, since each count is even. A two-critical-point height
gives a double of two smooth disks, which is standard because any
boundary diffeomorphism is isotopic to an orthogonal map
\cite{Hatcher1983Smale}. With four critical points the ranks in
the Morse complex give one extra pair in adjacent indices.
An extra \(0,1\) pair joins two distinct disks and cancels; duality
handles \(3,4\). For a middle pair reverse the function if
necessary to have one 2- and one 3-handle. The disk with those
two handles has \(H_2=0\), since only the final 4-handle is missing.
It is a standard disk by the handle consequence of Property R
\cite[Corollary 8.6]{Gabai1987FoliationsIII}. Equivalently the
2-handle has zero framing on a knot whose surgery is \(S^1\times S^2\);
Property R makes the knot trivial. Its trace is \(S^2\times D^2\); gluing it to the dual \(S^1\times D^3\) with any boundary marking gives \(S^4\) \cite[Theorem A]{LaudenbachPoenaru1972Handlebodies}. The final
disk gluing gives the same standard sphere. Hence every immersion
of a nonstandard candidate must have \(J\ge\omega\).

In fact every Morse function on a nonstandard candidate has at least
two index-two critical points. Reorder handles by index. The union
of its 0- and 1-handles is connected: no higher attaching region
can connect different components. Amalgamate the disks and 1-handles
of a spanning tree of this connected graph into one disk. The
remaining 1-handles reattach to its actual boundary. Do the same
for the disjoint dual 0- and 1-handlebody. Transport the intervening
framed 2-handle attachments with these boundary identifications.
Their number \(m_2\) is unchanged, and Euler characteristic now says
\(m_1+m_3=m_2\). If \(m_2\le1\), the preceding disk and single
middle-pair arguments give the standard sphere.

Let \(U\) be the open set where the shape operator has two eigenvalues
of each sign, counted with multiplicity. At a Morse critical point
the height Hessian is the second fundamental form up to sign. Reversing
the normal preserves index two. Thus, writing \(m_2(u)\) for the
original height count, the indexed area formula gives
\begin{equation}\label{eq:middle-mass}
 \frac1\omega\int_{S^4}m_2(u)\,du=\frac2\omega\int_U\det A\,dV_g .
\end{equation}
In particular the integral on the right before the factor \(2/\omega\)
is at least \(\omega\) on a nonstandard candidate. The comparison
of the Weyl tensor on middle signatures with the indexed height integral
is due in general dimension to Onti and Vlachos
\cite[\S\S2--3]{OntiVlachos2017Hypersurfaces}. Here the smooth
four-dimensional handle conclusion gives the lower count even though
the middle Betti number vanishes.

In a principal orthonormal frame let \(k_1,\ldots,k_4\) be the
principal curvatures. The Gauss equation gives the scalar curvature
\(R=3H^2/4-|A^\circ|^2\). On both chiral Weyl blocks its eigenvalues
are \((s_j-\bar s)/2\), where \(s_j=k_1 k_l+k_m k_n\) for the
three pairings and \(\bar s\) is their mean. Their squared Hilbert--Schmidt norm is
\(\frac1{12}\sum_{j<k}(s_j-s_k)^2\). On \(U\), write the curvatures
as \((-a,-b,c,d)\) with \(a,b,c,d>0\). The positive pair sum
\(ab+cd\) is at least \(2\sqrt{\det A}\) and each negative pair
sum is at most \(-2\sqrt{\det A}\). Therefore
\begin{equation}\label{eq:middle-weyl}
             |W^\pm|^2_{\rm HS}\ge\tfrac83\det A\quad\text{on }U .
\end{equation}
Equality requires \(ab=cd,\ ac=bd,\ ad=bc\), so \(a=b=c=d\).
Equality throughout a nonempty open subset of \(U\) is impossible.
Indeed the two smooth principal distributions there have rank two
and eigenvalues \(k,-k\) with \(k>0\). Codazzi, tested on orthogonal
vectors in either distribution, gives the directional derivative
of its eigenvalue zero along that distribution. Together they give
\(dk=0\). For \(X\) in the first distribution and \(Y\) in the
second, \((\nabla_X A)Y=(\nabla_Y A)X\) now lies in both
distributions and vanishes. Metric compatibility then shows each
distribution is parallel in all directions. A mixed sectional
curvature would be zero, contrary to the Gauss value \(-k^2\).
Thus when \(\int_U\det A>0\) the integrated comparison is strict.
Equations~\eqref{eq:middle-mass}--\eqref{eq:middle-weyl} force
\(E_+(g)>(8/3)\omega=64\pi^2/9\) on a nonstandard candidate.
More precisely the excess is the sum of three nonnegative terms,
\[
 E_+(g)-\tfrac83\omega =
 E_+(g;M\setminus U)
 +\int_U\bigl(|W^+|^2_{\rm HS}-\tfrac83\det A\bigr)\,dV_g
 +\tfrac43\omega\left(\frac1\omega\int_{S^4}m_2(u)\,du-2\right).
\]
If \(E_+(g)<8\pi^2\), the mean index-two count is less than \(9/4\).
All regular counts are at least two, so the proportion at least
three is less than \(1/4\). This proves the height assertion.
In particular four mutually orthogonal directions can be chosen
with this count simultaneously: the first four vectors of a random
orthonormal frame each have uniform marginal direction, and the
union of the four failure events has probability less than one.
After the two spanning-tree reductions such a two-middle-handle
description has \((m_1,m_3)=(0,2),(1,1)\) or \((2,0)\).

If the four trace-free principal curvatures are
\[
 (x+y+z,\ x-y-z,\ -x+y-z,\ -x-y+z)/2,
\]
then \(N=x^2+y^2+z^2=|A^\circ|^2\) and the three Weyl eigenvalues
are \((x^2-N/3)/2,(y^2-N/3)/2,(z^2-N/3)/2\).
Consequently the pointwise deficit satisfies the sharp inequality
\begin{equation}\label{eq:shape-deficit}
 D:=\tfrac16|A^\circ|^4-|W^\pm|^2_{\rm HS}
   =\tfrac12(x^2y^2+x^2z^2+y^2z^2)
   \ \ge\ \tfrac89(-\det A)_+ .
\end{equation}
For negative determinant, reverse normal if necessary, permute,
and rescale the principal curvatures by a positive constant to
write them as \((-s,1-s,1-s+v,1-s+v+u)\), with \(0<s<1\),
\(v,u\ge0\). Put \(z_0=1+2v+u\), \(b=1-2s\).
Then \(32D=2z_0^2(1+u^2)+(1-u^2)^2\ge2z_0^2+1\), whereas
\(16(-\det A)=(1-b^2)((z_0+b)^2-u^2)\). Cauchy--Schwarz gives
\[
 (1-b^2)(z_0+b)^2\le(2z_0^2+1)(1-b^2)(1/2+b^2)
                         \le\tfrac9{16}(2z_0^2+1).
\]
This proves \eqref{eq:shape-deficit}; \((-1,3,3,3)\) realizes
equality. For nonnegative determinant the formula for \(D\)
suffices. Integrating on a nonstandard candidate yields
\[
                E_+(g)\le B^2/6-64\pi^2/27 .
\]

For \(B\le(16\sqrt2/3)\pi\) the upper bound is \(64\pi^2/9\),
which contradicts the strict middle-index estimate. This proves
the third recognition criterion.

The same calculation gives a separate sign test for an actual
metric constructed using this representation. Use the sharp submanifold Sobolev inequality in
codimension two \cite[Corollary 2 and \S2]{Brendle2021Isoperimetric}
on \(F\) composed with \(\mathbb R^5\subset\mathbb R^6\).
The proof applies on the immersed domain: its contact function
takes a minimum there, and the area formula for the normal-bundle
map counts the preimages of each ambient point. The extra normal
has zero second fundamental form. Writing \(\kappa=|B^4|=\pi^2/2\),
the inequality for \(w=|f|^3\) and Cauchy--Schwarz give
\[
 4\kappa^{1/4}\|w\|_{4/3}\le\int_M\sqrt{|dw|^2+H^2w^2}\,dV_g,
 \qquad 16\sqrt\kappa\,\|f\|_4^2\le9\|df\|_2^2+\|Hf\|_2^2 .
\]
Thus the closed conformal-Laplacian form is at least
\((32\sqrt\kappa/3-B)\|f\|_4^2+\|Hf\|_2^2/12\).
At \(B=32\sqrt\kappa/3\), apply this to the positive first
eigenfunction. The mean curvature is nonzero somewhere since the
coordinate functions of a closed minimal immersion into Euclidean
space would be harmonic constants. This proves strict positivity
also at that endpoint. For a larger \(B\), one can test the smaller penalty
\(\|(|A^\circ|^2-H^2/12)_+\|_2\) against \(32\sqrt\kappa/3\)
with strict inequality. On a nonstandard candidate this would give
positive-class entry whenever also \(B<(4\sqrt{35}/3)\pi\), by
the integrated deficit estimate. All bounds refer to the same immersion.
\end{proof}

The immersion theorem supplies the first derivatives needed for
an immersion. The quantitative tests use its second derivatives
over the whole domain. Constructing an immersion satisfying either
test on an arbitrary smooth type remains an existence question.


\section{What would complete the argument}\label{sec:open}

For each smooth homotopy four-sphere left by the independent recognition
tests, it would suffice to construct a metric satisfying Proposition~\ref{prop:entry} and to exclude an attained minimum in \eqref{eq:band} on every
nonstandard type. The first assertion supplies a positive class to
which the fixed-type theorem applies. The second would then contradict
nonstandardness. The fixed-band finiteness theorem gives no uniform
distance of \(I^+(M)\) from the upper threshold as \(M\) varies.
Existence of the entry metric and rigidity of the attained minimum
are the two open geometric assertions in this reduction.

\section*{Provenance and assistance}
Edward B. Baker III developed this research with assistance from OpenAI
and Anthropic language models. The 26 November 2026 editorial revision
uses Claude Fable 5.2's review of the master at \texttt{c8fb08d}, the later
marked-surgery results, and the working proof of fixed-type attainment
\cite{Baker2026BourAttainment}. The companion proofs and their model
readings are research manuscripts. Human specialist verification of the
recognition gap and fixed-type attainment remains outstanding. Actual
serving model and reasoning effort for this revision were unavailable.

\bibliographystyle{unsrturl}
\bibliography{references}
\end{document}
